\section{Evaluation: what would establish broader value?}\label{sec:evaluation}

RRE needs evaluation at the level of mathematical progress and human use. More generated nodes, longer traces, or persuasive narratives are weak proxies. The framework should be judged by whether it helps produce correct, relevant, reusable results and whether others can understand and extend them at reasonable cost.

\subsection{Prospective research comparisons}

A first study of goal-directed exploration should use several predeclared mathematical objectives with explicit intermediate milestones. Problems should vary in their dependence on source retrieval, quantitative bounds, and transfer between settings. Include tasks on which a plausible route is false or a milestone is out of reach within the budget. Report all assigned tasks and stopped runs, including negative outcomes.

At least three conditions would be informative: assistance using ordinary research notes and document retrieval; assistance using a static curated reasoning graph; and RRE using the same initial graph with target generation and incremental updates. Match available sources, model versions, computational resources, time, and access to verification tools. Record the operating mode, human contributions, and any handoffs; goal-directed autonomous and collaborative RRE runs need not have the same labor requirements. Count the cost of constructing and curating the initial graph. Otherwise, an apparent improvement could reflect additional preparation rather than the recursive method.

Assessment should use reviewers with relevant expertise and, where possible, conceal the experimental condition. They should judge correctness, scope, novelty relative to documented prior work, and usefulness to the declared objective. Report review time, disagreements, and uncertainty. A result obtained only after substantial human repair should include that repair in both the outcome and resource accounting.

Useful endpoints include time to the first accepted intermediate result, fraction of proposed routes rejected for scope errors, number of previously documented failures repeated, and success of a later researcher using the recorded result. These measure different things and should not be combined into a single score without justification. The number of attempted subproblems is informative only with their difficulty and outcomes.

Open-ended exploration needs a separate comparison with matched starting collections and budgets. Reviewers should assess the mathematical value of the generated objectives as well as the correctness and reuse of their outcomes. Record rejected objectives and unsuccessful investigations; allowing an agent to choose easy questions must not be mistaken for improved discovery. Research-strategy improvement requires comparisons before and after a method change on held-out tasks, with errors, resource use, and human repair included. Neither capability is evaluated in the present case study.

\subsection{Ablations and cumulative effects}

To test the value of recursion, freeze the graph in one condition and allow reviewed increments in another. To test the value of negative knowledge, remove access to prior failed routes while retaining successful arguments. To test scope tracking, compare full records with otherwise identical records lacking explicit parameter and quantifier fields. Such ablations should be used in evaluation copies, not by corrupting an active research collection.

A cumulative knowledge claim requires a second stage. Freeze the outputs of one investigation, then ask a new researcher or a fresh system to solve a related target using that record. Measure whether the added lemma, counterexample, or quantitative bound is correctly retrieved and reused. The case study provides several connected developments within one investigation: finite certificates, a uniform pencil bound, balanced transfer, and the all-target two-term theorem. A subsequent-user study would test whether those additions also help a separate investigation.

\subsection{Review, replay, and correction experiments}

A replay study should give an independent evaluator the artifact package and a declared task: reconstruct the accepted argument, reproduce a certificate, identify a missing premise, or explain the stopping decision. Success requires more than opening the files. The evaluator must be able to locate the necessary source versions and distinguish what was proved from what was tested.

Correction propagation can be evaluated with controlled changes to a copied graph. For example, withdraw a premise used by two derivations while retaining an unaffected alternative proof. Check whether the system flags exactly the affected routes and preserves supported alternatives. Deliberately introduced scope errors can test whether uniform and fixed-parameter claims are distinguished. Automated structural checks and expert mathematical judgments should be reported separately.

\subsection{Human learning and collaboration}

Educational evaluation should compare a polished proof, the same proof with a static dependency graph, and a curated exploration history. Outcomes should include transfer to a new problem, identification of missing hypotheses, quality of proposed intermediate targets, and retention after a delay. Time on task and prior knowledge matter. A visually engaging display is not evidence of improved mathematical understanding.

Collaboration studies should ask whether contributors can resume one another's work, assign review to specialists, and preserve attribution under concurrent revisions. They should also measure the burden of maintaining the record. If recording every action consumes more time than it saves, RRE will need better selective capture and summarization. A useful system should support different levels of detail while making the recording policy explicit.

\begin{table}[htbp]
\centering\small
\begin{tabularx}{\linewidth}{@{}p{3cm}YY@{}}
\toprule
Claim to evaluate & Evidence available here & Further evidence needed \\
\midrule
Recorded recursive investigation is feasible & One run with successive generated targets and six graph revisions & Multiple independently executed runs \\
Mathematical outcomes are correct & Full theorem proofs, an explicit certificate, and exact replay & Independent mathematical review; selected formalization \\
Graphs improve discovery & No controlled comparison & Matched baselines and preregistered milestones \\
Records support cumulative reuse & Reuse within the selected run & Later users and held-out objectives \\
Replay reduces review effort & Saved artifacts and internal replay & Independent reconstruction and cost measurement \\
Exploration paths support learning & Proposed exercises and interfaces & Student studies with transfer outcomes \\
\bottomrule
\end{tabularx}
\caption{Current evidence and the evaluations required for stronger claims.}
\end{table}
